\documentclass[11pt,a4paper,fontset=fandol,scheme=plain]{ctexart}

\usepackage[margin=25mm]{geometry}
\usepackage{microtype}
\usepackage{xcolor}
\usepackage{hyperref}

\definecolor{accent}{HTML}{285FAE}
\hypersetup{
    unicode=true,
    pdftitle={理性是否真正理性：一个概率视角的讨论},
    pdfauthor={Xayah Hina},
    colorlinks=true,
    linkcolor=accent,
    urlcolor=accent
}

\setlength{\parindent}{2em}
\setlength{\parskip}{0pt}
\setcounter{secnumdepth}{3}

\title{理性是否真正理性：一个概率视角的讨论}
\author{Xayah Hina}
\date{2026-07-15}

\begin{document}
    \maketitle

    \begin{abstract}
        本文目标：探讨理性介入和规划执行在概率视角下对于现实优化的实际效果问题。
    \end{abstract}


    \section{Motivation}
    最近几年，尤其是LLM的实证成功让我对于概率这个领域的思考越来越多，体会也会来越频繁。近期对于


    \section{More is Different? 对世界概率性的思考历程}
    这一章节回顾一下自己从本科那会开始对于整个世界概率性（注：此处的概率学并非完全等价于数学中的概率论）的整个思考历程。
    
    \subsection{从大数定律和中心极限定律说起}

    对于世界的概率性的初步认识自然是来自于修概率论与数理统计这门课的时候。说实话当时对于前几章的排列组合求概率等中学已熟知的内容并无太多兴趣，直到第五章（对于这一章我印象非常清晰，当时还特意上课时多看了很多遍）老师讲大数定律/中心极限定理时，突然眼前一亮大脑清澈，这种感受我直接仍记忆犹新。尤其是当时老师提到去一个面包店排队的人从微观上是随机的但是一旦从一个月来统计就会发现，一天中人群排队的按每小时统计的值会回归到一个\textbf{稳定的值}。
    这个现象说实话挺“显而易见”的但是细细思考下去就会触及到我之前从没有想过的一面。之前我总是以为随机就是随机，掷骰子是随机，分子随机运动是随机，怎样的人什么时候去网购一键什么样的东西也是随机。那时我就在想，这么说其实这种“随机其实并不是随机”而是一种确定意义下的随机？或者说完全可以视作一种隐藏的确定值？这个\textbg{隐藏}给了当时的我很大启发。果不其然下一节课就讲到了中心极限定律，这个世界上的万事万物都和正态分布有着千丝万缕的联系（难怪要叫Normal Distribution）。
    在奶茶店一个月每天上午10点排队的人数就是近似正态分布。

\end{document}
